%\section{Systematic Uncertainties} 

\begingroup
\squeezetable
\begin{table}
\begin{tabular}{cccccccc}
$Q^2_{QE}$ (GeV$^2$) & I & II & III & IV & V & VI & Total \\
\hline
$0.0 - 0.025$ & 0.05 & 0.04 & 0.00 & 0.02 & 0.11 & 0.02 & 0.13 \\
$0.025 - 0.05$ & 0.05 & 0.04 & 0.01 & 0.01 & 0.11 & 0.02 & 0.13 \\
$0.05 - 0.1$ & 0.05 & 0.04 & 0.01 & 0.01 & 0.11 & 0.01 & 0.13 \\
$0.1 - 0.2$ & 0.04 & 0.04 & 0.01 & 0.01 & 0.11 & 0.01 & 0.12 \\
$0.2 - 0.4$ & 0.03 & 0.06 & 0.01 & 0.02 & 0.11 & 0.01 & 0.13 \\
$0.4 - 0.8$ & 0.05 & 0.07 & 0.02 & 0.03 & 0.11 & 0.01 & 0.15 \\
$0.8 - 1.2$ & 0.11 & 0.11 & 0.02 & 0.02 & 0.11 & 0.02 & 0.20 \\
$1.2 - 2.0$ & 0.13 & 0.15 & 0.04 & 0.04 & 0.12 & 0.02 & 0.23 \\
\hline
\end{tabular}
\caption{Fractional systematic uncertainties on $d\sigma/dQ^2_{QE}$ associated with muon reconstruction (I), recoil reconstruction (II), neutrino interaction models (III), final state interactions (IV), flux (V) and other sources (VI).  The final column shows the total fractional systematic uncertainty due to all sources.}
\label{tab:systematics}
\end{table}
\endgroup


The main sources of systematic uncertainty in the differential
cross-section measurement are due to: the reconstruction of the muon;  
the reconstruction and detector response for hadrons; 
 the neutrino interaction model; final state interactions; 
and the neutrino flux.  These uncertainties are evaluated by 
repeating the cross-section analysis with systematic shifts 
applied to the simulation and their effect is shown in 
Tab.~\ref{tab:systematics}.

%The second largest uncertainty comes from the muon momentum scale in MINOS. 
Uncertainties in the muon energy scale have a direct impact on $Q^{2}_{QE}$ and result in a bin migration of events. The bulk of the uncertainty comes from \minos, which reconstructs muon energy by range (for stopping tracks) and curvature (exiting). There is a 2.0\% uncertainty in the range measurement, due to the material assay and imperfect knowledge of muon energy loss\cite{Michael:2008bc}. Using muon tracks which stop in the detector, we compare the momentum measured by range and curvature to establish an additional uncertainty of 0.6\% (2.6\%) on curvature measurements above (below) \unit[1]{GeV/c}.
We also account for subdominant uncertainties on the  
 energy loss in \minerva, systematic offsets in the 
beam angle, mis-modeling of the angular and position resolution, 
and tracking efficiencies.

The systematic error on the recoil energy measurement is due to the
uncertainty in the \minerva detector energy scale set by muons and
differences between the simulated calorimetric response 
to single hadrons and the response measured by the test beam program.
Additional uncertainties are due to differences between the Geant model of
neutron interactions and thin target data on neutron scattering in
carbon, iron and copper\cite{Abfalterer:2001gw,Schimmerling:1973bb,Voss:1956,Slypen:1995fm,Franz:1989cf,Tippawan:2008xk,Bevilacqua:2013rfq,Zanelli:1981zz}.   We evaluate further sources of systematic error 
by loosening analysis cuts on energy near the vertex 
and on extra isolated energy depositions, repeating the 
fit to the background and subsequent analysis, and assigning
 an uncertainty to cover the difference.  

Predictions for $Q^{2}_{QE}$ and recoil energy distributions for neutrino-induced background processes are based upon the GENIE generator. We evaluate the systematic error by varying the underlying model tuning parameters according to their uncertainties~\cite{Andreopoulos201087}. These include parameters governing inelastic interactions of neutrinos with nucleons and those that vary the final state interactions.  

The systematic error on the antineutrino flux arises from uncertainties in hadron production in the NuMI target and beamline, and from imperfect modeling of the beamline focusing and geometry~\cite{Pavlovic:2008zz}.  Where hadron production is constrained by NA49 data\cite{Alt:2006fr}, the NA49 measurement uncertainties dominate.  The uncertainty on other interactions is evaluated from the spread between different Geant4 hadron production models~
%\cite{Agostinelli2003250}. 
\cite{Agostinelli2003250,1610988}
The absolute flux uncertainties are large, with a significant $E_\nu$ dependence, but mostly cancel in a measurement of the shape of $d\sigma /dQ^2_{QE}$.




%The largest uncertainties in this measurement are due to uncertainties in the %flux, which themselves are dominated by those associated with 
%hadron production.  When hadron production is
%constrained by the NA49 data, the NA49 measurement uncertainties are
%the dominant source~\cite{Alt:2006fr}.  For interactions resulting from
%pions outside the NA49 acceptance or which result from
%reinteractions of pions or secondary nucleons, the uncertainties are
%evaluated by taking the differences between different Geant4 hadron
%production models: FTFP (BERT), QGSP (BERT)
%and FLUKA.  The flux uncertainty from the beamline geometry and
%focusing fields is dominated by an imperfect knowlege of the horn current and
%radial distribution of that current in the horn's inner
%conductor~\cite{Pavlovic:2008zz}.  The flux uncertainties, although dominant at
%almost all $Q^2_{QE}$ as
%shown in Fig.~\ref{fig:syst_err}, are nearly constant with
%$Q^2_{QE}$ and are negligible in a
%measurement of the shape of the $Q^2_{QE}$ distribution.

%The second largest uncertainty in the differential cross-section
%comes from the uncertainty in muon momentum scale in the MINOS
%spectrometer.  Uncertainties in modeling of muon energy loss
%processes, in tracking near the track endpoint, 
%and limitations in a test beam calibration limit the accuracy of the
%momentum measured by range of the muons in the MINOS spectrometer to
%2.0\%.  The calibration of momentum measured by track curvature studies %differences between momenta measured by range to that measured by
%curvature, and comparisons of these differences in data and simulation
%lead to an additional uncertainty of 0.6\% (2.6\%) for whose curvature momentum %is above (below) 1~GeV; this uncertainty is uncorrelated with that of the range
%measurement.  


%The fit of the
%recoil distribution to that measured in data to determine the
%background fraction reduces these uncertainties in the final result.

%\begin{figure}[tp]
%\centering
%\includegraphics[width=75mm, trim=2cm 0 2cm 0]{figures/err_summary_total_bayes_variedbackgrounds_fit_sometimesUnsmeared_standard_minerva}
%\caption{Systematic uncertainty summary for the 
%cross-section shape measurement as a function of $Q^2_{QE}$.  }
%\label{fig:syst_err}
%\end{figure}
